% ===========================================================================
\chapter{\label{ap:prompts}Prompts de geração}

\section{Condição P: geração a partir do fragmento-fonte}

Prompt usado na versão avaliada para gerar uma questão. O marcador \texttt{\{conteudo\}} é substituído pelo fragmento-fonte, truncado em 6.000 caracteres. Fonte: \texttt{backend/app/llm.py} do protótipo.

\begin{small}
\begin{verbatim}
Você é um professor criando uma questão de estudo a partir de um
material didático.
Com base APENAS no trecho abaixo, crie 1 questão de múltipla escolha
em português que avalie a compreensão do conteúdo.

Requisitos:
- Enunciado claro, contextualizado e respondível somente com o
  conteúdo do trecho.
- Exatamente 4 alternativas plausíveis, sendo apenas 1 correta.
- Evite alternativas obviamente erradas ou sem sentido.

TRECHO:
{conteudo}

Responda SOMENTE com um JSON válido no formato:
{"pergunta": "...", "opcoes": ["...", "...", "...", "..."],
 "correta": 0, "justificativa": "..."}
\end{verbatim}
\end{small}

Na sessão de estudo, as questões são geradas em lotes com o prompt a seguir, em que \texttt{\{trechos\}} contém até cinco fragmentos numerados, cada um truncado em 1.500 caracteres, e \texttt{\{n\}} é o número de fragmentos.

\begin{small}
\begin{verbatim}
Você é um professor criando questões de estudo a partir de um
material didático.
Para CADA trecho numerado abaixo, crie 1 questão de múltipla escolha
em português (4 alternativas plausíveis, exatamente 1 correta),
respondível apenas com o conteúdo daquele trecho.

TRECHOS:
{trechos}

Responda SOMENTE com um JSON válido no formato, com exatamente {n}
itens na ordem dos trechos:
{"questoes": [{"pergunta": "...", "opcoes": ["...", "...", "...",
 "..."], "correta": 0, "justificativa": "..."}]}
\end{verbatim}
\end{small}

\section{Condições P+V e S}

Na condição P+V, o prompt da condição P recebe, após o trecho, a seção a seguir, com os dois fragmentos mais próximos recuperados do banco vetorial; a instrução de ancoragem no trecho principal é mantida.

\begin{small}
\begin{verbatim}
CONTEXTO ADICIONAL (trechos relacionados do mesmo material; use
apenas para compreender o trecho principal, a questão deve ser
respondível com o TRECHO):
{vizinhos}
\end{verbatim}
\end{small}

Na condição S, o bloco \texttt{TRECHO} é substituído pelas linhas abaixo, e a frase ``Com base APENAS no trecho abaixo'' é substituída por ``Sobre o tema abaixo''. \definir{validar a redação final das condições P+V e S antes da coleta e congelá-la no repositório}.

\begin{small}
\begin{verbatim}
MATERIAL: {titulo_do_material}
TEMA: {tema_do_fragmento}
\end{verbatim}
\end{small}

% ===========================================================================
\chapter{\label{ap:rubrica}Rubrica de inspeção das questões}

A Tabela~\ref{tab:rubrica} apresenta a ficha usada por cada avaliador. O avaliador vê o fragmento-fonte, o enunciado, as quatro alternativas, a indicação da correta e a justificativa, sem a identificação da condição.

\begin{table}[htbp]
\centering
\caption{Ficha de avaliação de cada questão, baseada nas diretrizes de Haladyna et al. e na taxonomia revisada de Bloom \cite{AK01}. Fonte: elaborada pelo autor.}
\label{tab:rubrica}
\small
\begin{tabularx}{\textwidth}{l >{\raggedright\arraybackslash}X >{\raggedright\arraybackslash}p{3.8cm}}
\toprule
Critério & Pergunta ao avaliador & Opções \\
\midrule
C1 & A alternativa indicada como correta é sustentada pelo fragmento-fonte, sem exigir conhecimento externo? & Sim / Parcialmente / Não \\
C2 & A alternativa indicada como correta está correta, e nenhuma outra alternativa também está? & Sim / Não \\
C3 & Quantos distratores seriam plausíveis para um estudante que não domina o conteúdo? & 0 / 1 / 2 / 3 \\
C4 & O enunciado é inequívoco, gramaticalmente correto e compreensível sem ler as alternativas? & Sim / Não \\
C5 & Qual o processo cognitivo principal exigido? & Lembrar / Compreender / Aplicar ou superior \\
Obs. & Defeito observado, se houver (texto livre). & -- \\
\bottomrule
\end{tabularx}
\end{table}

Antes da avaliação principal, os dois avaliadores classificam juntos dez questões de exemplo, geradas de um material fora do corpus, e discutem as divergências, para alinhar a interpretação dos critérios. Essas questões não entram nos resultados.

% ===========================================================================
\chapter{\label{ap:usuarios}Desenho do estudo com participantes (trabalho futuro)}

Este apêndice descreve o estudo com participantes que dá continuidade a este trabalho, condicionado à aprovação do CEP-PUCRS. Ele incorpora as críticas feitas à versão anterior da proposta e parte dos resultados técnicos obtidos aqui.

\section{Pergunta e desenho}

A pergunta é se intercalar questões geradas pelo \ferramenta durante a leitura melhora a retenção do conteúdo, em comparação com a leitura do mesmo tipo de material sem questões. A comparação entre as duas condições se justifica porque a pergunta é sobre o efeito da intervenção; um desenho com uma única condição só permitiria estudar correlações dentro dela.

O desenho é intra-sujeitos: cada participante realiza as duas condições, com materiais distintos e pareados em extensão e dificuldade, em ordem contrabalanceada e com pelo menos 24 horas de intervalo \cite{WRH12}. A variável dependente principal é o desempenho em um pós-teste de retenção aplicado uma semana após cada sessão, composto por questões escritas por um especialista, independentes das geradas pela ferramenta, para que o teste não favoreça a condição com questões. As variáveis secundárias são o \fs médio e os eventos de celular, apenas em configurações que a QP2 mostrou viáveis.

\section{Participantes e análise}

Estudantes voluntários maiores de 18 anos, com \estr{notebook} com câmera. O número de participantes deve ser definido por análise de poder a partir de um tamanho de efeito de referência \referencia{tamanho de efeito do efeito de testagem ou de perguntas intercaladas em estudos com materiais comparáveis, para cálculo de poder}. A comparação entre condições usa o teste de Wilcoxon para amostras pareadas, adequado a amostras pequenas sem pressuposto de normalidade \cite{Fie13}, com tamanho de efeito e intervalo de confiança. A relação entre \fs e taxa de acertos na condição com questões é explorada pela correlação de Spearman com intervalo de confiança por \estr{bootstrap}.

\section{Instrumentos complementares}

Ao final, cada participante responde à \estr{System Usability Scale} (SUS) \cite{Bro96}, interpretada segundo Bangor et al.\ \cite{BKM08}, e a três perguntas de benefício percebido: se a ferramenta contribuiu para o seu aprendizado, se as questões ajudaram a perceber o que não havia entendido e se usaria a ferramenta para estudar para uma prova. As respostas abertas são analisadas por análise de conteúdo \cite{Bar11}, com codificação pelo autor, revisão pelo orientador e registro das divergências e de sua resolução.
